\section{Explicit residues and effective Brauer--Siegel}
\label{sec:normal}

Throughout, $D_K=|\operatorname{disc}(K)|$, $n_K=[K:\mathbf Q]$,
$(r_1,r_2)$ is the signature, $w_K=|\mu(K)|$, and
$\kappa_K=\operatorname*{Res}_{s=1}\zeta_K(s)$.
The regulator uses the usual logarithmic embedding with a factor $2$
at complex places, and is $1$ in unit rank zero. Thus
\begin{equation}\label{eq:classformula}
 h_KR_K=\frac{w_K\sqrt{D_K}}{2^{r_1}(2\pi)^{r_2}}\kappa_K.
\end{equation}
This is the analytic class number formula, in precisely the normalization
of Garcia--Lee \cite[(1), p.~2]{GarciaLee}.

\subsection{The constants in Stark's reduction}

Here are the inputs, with their original locators checked against the
1974 facsimile. Stark \cite[Lemma~3, pp.~139--140]{Stark} gives at most one
zero in
\[
 \operatorname{Re}s\geq1-\frac1{4\log D_K},\qquad
 |\operatorname{Im}s|\leq\frac1{4\log D_K};
\]
such a zero is real and simple. Write it as $\beta_0$ if it exists, and
set $\beta_1=\beta_0$ in that case, and
$\beta_1=1-(4\log D_K)^{-1}$ otherwise. His Lemma~4, equation~(11),
pp.~140--141, gives
\begin{equation}\label{eq:stark-residue}
 \kappa_K>C_4^{-1}(1-\beta_1),\qquad
 C_4=2\exp\left(\frac{21}{8}
                 +\frac{\gamma+2\log2}{4\log3}\right)<43.2.
\end{equation}
The displayed constant follows by taking $c_1=0$, $c_2=2/\log3$ in
Stark's equation~(7), p.~139, and $c_4=2c_3$ on p.~141.
Garcia--Lee \cite[\S3, pp.~16--17]{GarciaLee} independently records this
same extraction. In particular $4C_4<173$.
For an exact numerical check of the last strict inequality, one may use
$\gamma<289/500$, $\log2<347/500$, and $\log3>549/500$;
the exponent in $C_4/2$ is then less than $3073/1000$.
The rational Taylor-series certificate supplied with this paper proves
$8\exp(3073/1000)<173$ without floating-point comparisons.

Stark's Lemma~8, p.~145, has exactly the constants $4$ for a normal
extension of $\mathbf Q$ and $4n_K!$ otherwise: a real zero in the
corresponding interval descends to a quadratic subfield. Theorem~$1'$,
p.~148, equations~(24)--(25), is the resulting residue statement;
its no-quadratic-subfield clause uses the very same zero exclusion.
We extend that clause using Alp\"oge's theorem, not the other branch of
Theorem~$1'$ involving the constant $c_8$. No numerical assertion about
Stark's tentative value $c_8=\pi/6$ is needed here.

\begin{theorem}[NX01: normal fields, with an absolute constant]
\label{NX01}
Assume the supplied theorem of Alp\"oge in the sense fixed in the
introduction. For every normal number field $K/\mathbf Q$ of degree
$n_K\geq2$,
\begin{equation}\label{eq:normal-lower}
 \boxed{\quad
 \kappa_K>\frac1{173\log D_K},\qquad
 h_KR_K>
 \frac{w_K\sqrt{D_K}}
 {173\,2^{r_1}(2\pi)^{r_2}\log D_K}.
 \quad}
\end{equation}
More precisely, $173$ may be replaced by $4C_4$ with $C_4$ as in
\eqref{eq:stark-residue}. The field $\mathbf Q$ is handled separately:
$D=1$, $\kappa=hR=1$.
\end{theorem}

\begin{proof}
Suppose that a real zero $\beta$ lies in
$[1-(4\log D_K)^{-1},1)$. Since $D_K\geq3$, this interval lies strictly
to the right of $1/2$. Stark's Lemma~8 gives a quadratic subfield
$F\subset K$ with $\zeta_F(\beta)=0$. The factorization
$\zeta_F(s)=\zeta(s)L(s,\chi_{D_F})$ contradicts the supplied
Riemann and Dirichlet nonvanishing theorems at $\beta$.
Thus $\beta_0$ does not exist, and \eqref{eq:stark-residue} gives
$\kappa_K>(4C_4\log D_K)^{-1}>(173\log D_K)^{-1}$.
Equation~\eqref{eq:classformula} gives the second inequality.
The conclusion uses only zero exclusion near $1$ for normal Dedekind
zeta functions; it asserts no full GRH for them.
\end{proof}

\begin{theorem}[NX02: transfer through the Galois closure]
\label{NX02}
For every number field $K$ of degree $n\geq2$, under the same input,
\begin{equation}\label{eq:all-fields-lower}
 \boxed{\quad
 \kappa_K>\frac1{173n!\log D_K},\qquad
 h_KR_K>
 \frac{w_K\sqrt{D_K}}
 {173n!\,2^{r_1}(2\pi)^{r_2}\log D_K}.
 \quad}
\end{equation}
The factor $n!$ can be replaced by $1$ for normal $K$. In particular,
for a CM field $E$ of degree $2g$,
\begin{equation}\label{eq:cm-hR}
 h_ER_E>\frac{w_E\sqrt{D_E}}
 {173(2g)!(2\pi)^g\log D_E}.
\end{equation}
\end{theorem}

\begin{proof}
For clarity we spell out the normal-closure mechanism behind Stark's
Lemma~8, rather than treating a lower bound for the residue of the
closure as a lower bound for the residue of its subfield.
Let $L$ be the Galois closure, $m=[L:\mathbf Q]\leq n!$, and express $L$
as the compositum of $a\leq n$ conjugates of $K$. Stark's Lemma~7,
p.~143, gives
\[
 D_L\mid D_K^{am/n},\qquad \log D_L\leq m\log D_K\leq n!\log D_K.
\]
The Aramata--Brauer divisibility used in the proof of Stark's
Lemma~8 says that $\zeta_L/\zeta_K$ is entire, since $L/K$ is normal.
Consequently a real zero of $\zeta_K$ in
$[1-(4n!\log D_K)^{-1},1)$ would also be a zero of $\zeta_L$ in its
Stark interval, contrary to NX01's zero-exclusion argument.

If the exceptional zero $\beta_0$ in Stark's Lemma~4 for $K$ exists,
it therefore satisfies $1-\beta_0>(4n!\log D_K)^{-1}$; if it does not,
$1-\beta_1=(4\log D_K)^{-1}$ is at least this large.
Apply \eqref{eq:stark-residue} and then
\eqref{eq:classformula}. For a CM field $r_1=0$, $r_2=g$.
\end{proof}

\begin{theorem}[NX03: a quantitative Brauer--Siegel statement]
\label{NX03}
Let $K/\mathbf Q$ be normal, $n\geq2$, and put $L=\log D_K$. Then
\begin{align}
 \frac{w_K\sqrt{D_K}}{173\,2^{r_1}(2\pi)^{r_2}L}
 &<h_KR_K\notag\\
 &\leq\frac{w_K\sqrt{D_K}}{2^{r_1}(2\pi)^{r_2}}
       \left(\frac{eL}{2(n-1)}\right)^{n-1}.
 \label{eq:effective-BS}
\end{align}
In particular the following error interval is explicit:
\begin{align*}
 -\frac n2\log(2\pi)-\log(173L)+\log2
 &<\log(h_KR_K)-\frac12L\\
 &\leq\log(2n^2)+(n-1)\log\left(\frac{eL}{2(n-1)}\right).
\end{align*}
Along any sequence of normal fields with $D_K\to\infty$ and
$n/\log D_K\to0$, these inequalities give
$\log(h_KR_K)/\log\sqrt{D_K}\to1$ with the displayed effective error.
For arbitrary fields of fixed degree, the same conclusion and error
hold with the additional term $-\log(n!)$ in the lower bound.
\end{theorem}

\begin{proof}
The lower bound is NX01. The upper bound is Louboutin
\cite[Theorem~1(1), p.~1195]{NC-Louboutin2001}, inserted into the class number
formula. For the error interval use $w_K\geq2$ and
$2^{r_1}(2\pi)^{r_2}\leq(2\pi)^{n/2}$ in the lower estimate.
For the upper estimate, $w_K\leq2n^2$: indeed
$\varphi(w_K)\leq n$ and the prime-power product formula gives
$\varphi(t)^2\geq t/2$ for every positive integer $t$.
Divide the resulting logarithmic bounds by $L$. If $x=(n-1)/L\to0$,
then $x\log(e/(2x))\to0$, proving the assertion, and the other error
terms visibly tend to zero. The fixed-degree assertion uses NX02.
\end{proof}

\subsection{An independent upper bound relevant to the earlier programme}

\begin{proposition}[NX04: a class-number upper bound needs no new RH]
\label{NX04}
For any number field $K$ of degree $n\geq2$,
\begin{equation}\label{eq:minkowski-class-upper}
 h_K\leq\sqrt{D_K}\left(1+\tfrac12\log D_K\right)^{n-1}.
\end{equation}
Thus a previously established sufficient power
$2h_K|D_p|/|\mu(K)|$ can be bounded above effectively by substituting
\eqref{eq:minkowski-class-upper}, whenever the other factors of that
earlier statement have their specified finite meaning.
\end{proposition}

\begin{proof}
Minkowski's ideal-class bound \cite[Theorem~4.3, p.~70]{MilneANT} gives an integral ideal in every class
of norm at most
$M=(4/\pi)^{r_2}(n!/n^n)\sqrt{D_K}\leq\sqrt{D_K}$.
For the last inequality one may use $n!/n^n\leq2^{1-n}$ and
$r_2\leq n/2$; for $n\geq2$ the resulting factor
$2\pi^{-n/2}$ is at most $1$.
The Euler factors of $\zeta_K$ show that its nonnegative Dirichlet
coefficients are bounded coefficientwise by those of $\zeta(s)^n$.
If $d_n(m)$ denotes the latter coefficients, induction on $n$ gives,
for $X\geq1$,
\[
 \sum_{m\leq X}d_n(m)\leq X(1+\log X)^{n-1}.
\]
Indeed the induction step follows by summing
$\sum_{a\leq X}(X/a)(1+\log X)^{n-2}$ and using
$\sum_{a\leq X}1/a\leq1+\log X$.
Count ideals up to $\sqrt{D_K}$ to conclude.
\end{proof}

\paragraph{Geometric and programme interpretation.}
NX01--NX03 concern $\zeta_{\operatorname{Spec}\mathcal O_K}$ and the
arithmetic special-value combination $h_KR_K$; they give quantitative
information about the integral Picard group and the archimedean unit
lattice. They do not produce a decomposition of the corresponding
\'{e}tale complex, nor an integral $K$-theory calculation from zero
location. This distinction is exactly that of ZETA\_TWO ZT08, SV01--SV02.
The theorem extends ZETA\_TWO IA05's zero-exclusion reduction to explicit
residue bounds. NX04 applies to the sufficient-power expression in
SA06 B.4e as recorded in \texttt{INTAKE\_RESEARCH\_SA06\_20261120.md},
\S4.1; it does not repair the different refuted assertion B.4c.
A lower bound for $hR$ would require an \emph{upper} bound for $R$ to
give a lower bound for $h$. The CM argument below avoids that problem
by an exact regulator quotient. An upper bound for $hR$, in contrast,
needs a regulator lower bound to give an upper bound for $h$.

No manuscript statement of these new explicit estimates has been
identified in Grothendieck's archive. The exact programme anchors
available in the supplied record are SA06 B.4a--e and the determinant
discussion of cote~60, scans~3, 5, 7, 9, as recorded in ZETA\_TWO AR01
and its archive section. These anchors are cited as inherited readings,
not as new facsimile inspections or claims that the archive contains
the present analytic theorem.
